\documentclass{article}
\usepackage{amsthm}
\usepackage{amssymb}
\usepackage{amsmath}

\newtheorem{mydef}{Definici\'on}
\newtheorem{prop}{Proposici\'on}

\begin{document}

\title{Variables aleatorias discretas}
\author{nsm}

\maketitle


\begin{mydef}[v.a. cont\'inua]
  Una v.a. es \emph{cont\'inua} si existe una funci\'on
  \[
    f:\mathbb{R} \to \mathbb{R}_{\geq 0}
  \]
  tal que
  \[
    P(X \in A) = \int_A f(x) \ dx , \quad \forall A \subseteq \mathbb{R}
  \]
\end{mydef}

\begin{mydef}[Funci\'on de densidad]
  Dicha funci\'in se denomina \emph{funci\'on de densidad}.x
\end{mydef}

\paragraph{Propiedades de la funci\'on de densidad}
  Para que una funci\'on \(f(x)\) sea una funci\'on de densidad
  debe cumplir:
\begin{enumerate}
    \item \( f(x) \geq 0, \quad \forall x \)
    \item \( \int_{-\infty}^{+\infty} f(x) \ dx = 1\)
\end{enumerate}
\

\begin{mydef}[Funci\'on de distribuci\'on acumulada]
  La \emph{distribuci\'on} de una v.a. cont\'inua \(X\) con
  densidad \(f(x)\) se define para todo $x \in \mathbb{R}$ como:

  \[
    F_X(x) = \int_{-\infty}^{+\infty} f(t) \ dt
  \]
\end{mydef}

\paragraph{Propiedades de la funci\'on de distribuci\'on acumulada.}
Sea \(X\) una v.a. conr\'inua, entonces:
\begin{enumerate}
    \item $ \forall x \in \mathbb{R} \quad F_X(x) \in [0,1] $
    \item $ \forall x, y \quad x < y \implies F_X(x) \leq F_X(y) $
    \item \(F_X(x)\) es cont\'inua en todo punto.
    \item \(\lim_{x\to +\infty} F_X(x)=1 \land
      \lim_{x\to -\infty} F_X(x) = 0\)
\end{enumerate}
\

\begin{prop}
  Sean $a, b \in \mathbb{R} $ y $a \leq b$ entonces
  \[
    P(a \leq X \leq b) = F_X(b) - F_X(a)
  \]
\end{prop}

\begin{mydef}[Esperanza]
\end{mydef}

\paragraph{Propiedades de la Esperanza}
\begin{enumerate}
\item \(E[aX+b] \quad = \quad aE[X]+b \)
\item Si $P(X = c) = 1$, entonces $E[X]= c$
\end{enumerate}

\begin{mydef}[Varianza]
  Sea $X$ una v.a. discreta con f.p.p. $p_X(x)$ y esperanza $\mu_X$.
  La \emph{varianza}  de $X$ es:
  \[
    V[X] = E[(X - \mu_X)^2]
  \]
\end{mydef}

\begin{prop}[F\'ormula para la varianza]
  \[
    V[X] = E[X^2] - E[X]^2
  \]
\end{prop}
\end{document}
